\documentclass[11pt,a4paper]{article}

\usepackage{amsmath,amssymb,amsthm}
\usepackage{booktabs}
\usepackage{array}
\usepackage[colorlinks=true,linkcolor=blue,citecolor=blue,urlcolor=blue]{hyperref}
\usepackage{geometry}
\geometry{margin=2.5cm}
\usepackage{microtype}

% ---- notation ----------------------------------------------------------
\newcommand{\Es}{E^{*}}
\newcommand{\Order}{\mathcal{O}}
\newcommand{\R}{\mathbb{R}}
\newcommand{\Z}{\mathbb{Z}}
\newcommand{\norm}[1]{\left\lVert #1 \right\rVert}
\newcommand{\abs}[1]{\left\lvert #1 \right\rvert}
\newcommand{\xx}{\mathbf{x}}
\newcommand{\yy}{\mathbf{y}}
\newcommand{\pp}{\mathbf{p}}
\newcommand{\dd}{\mathbf{d}}
\newcommand{\Htwo}{\mathcal{H}^2}
\newcommand{\Hmat}{\mathcal{H}}
\newcommand{\code}[1]{\texttt{#1}}

% ---- theorem-like environments ----------------------------------------
\newtheorem{theorem}{Theorem}
\newtheorem{lemma}{Lemma}
\newtheorem{proposition}{Proposition}
\newtheorem{corollary}{Corollary}
\theoremstyle{definition}
\newtheorem{definition}{Definition}
\theoremstyle{remark}
\newtheorem{remark}{Remark}

\title{The black-box FMM ($\Htwo$) operator in Hcontact:\\[4pt]
a detailed derivation, design rationale, and validation}
\author{Hcontact --- theory companion to \texttt{doc/theory/h2\_fmm.tex}}
\date{2026-06-30}

\begin{document}
\maketitle

\begin{abstract}
This document is the in-depth companion to the concise note
\texttt{doc/theory/h2\_fmm.tex}. It explains, from first principles and at
textbook depth, the matrix-free hierarchical operator (\code{backend="h2"})
that Hcontact uses to apply the discrete Boussinesq half-space flexibility
matrix $S$ without ever forming it. The operator is a black-box fast multipole
method (bbFMM) in the sense of Fong and Darve~\cite{FongDarve2009}, built on
tensor-product Chebyshev interpolation and specialised to the
translation-invariant kernel on a uniform $N_s\times N_s$ grid. We derive every
formula directly from the source files \code{boussinesq\_kernel.\{hpp,cpp\}},
\code{cheb\_basis.\{hpp,cpp\}}, \code{uniform\_quadtree.\{hpp,cpp\}}, and
\code{h2\_operator.\{hpp,cpp\}}, citing line numbers where the correspondence is
load-bearing. We cover: the physical contact problem and its Green's function;
the Love~(1929) cell integral and the lookup table; why the dense and classical
$\Hmat$/ACA representations are wasteful for this kernel; the uniform quadtree
and a proof that the near field and the per-level interaction lists tile every
pair of unknowns exactly once; degenerate-kernel and Chebyshev interpolation
theory with convergence bounds; the full derivation of the five transfer
operators P2M, M2M, M2L, L2L, L2P, with the small-matrix forms exactly as
coded; the translation-invariant caching that collapses storage to $\Order(N)$;
operation and storage counts; accuracy analysis; a justification of every
default parameter; validation against the dense operator, the classical
$\Hmat$-matrix, and the FFT solver Tamaas; and limitations.
\end{abstract}

\tableofcontents
\bigskip
\hrule

\section{Introduction and reading guide}
\label{sec:intro}

Hcontact solves normal elastic contact between a rigid rough indenter and an
elastic half-space. The numerical bottleneck is the repeated application of a
dense, fully populated influence matrix $S\in\R^{N\times N}$, $N=N_s^2$, inside
a projected conjugate-gradient (PCG) iteration. Forming $S$ explicitly is
quadratic in storage and matrix--vector work, which becomes impossible well
before the grid sizes of interest ($N_s$ in the thousands, $N\sim10^{6}$--$10^{8}$).
This document describes the operator that removes that bottleneck: a
matrix-free $\Htwo$/FMM apply with $\Order(N)$ memory and $\Order(N)$
matrix--vector cost.

The presentation is deliberately self-contained and proceeds in the logical
order of the algorithm. Sections~\ref{sec:physics}--\ref{sec:cost} set up the
physics, the discretisation, and the cost problem that motivates a hierarchical
method. Sections~\ref{sec:tree}--\ref{sec:bbfmm} develop the geometry, the
approximation theory, and the five transfer operators. Section~\ref{sec:cache}
explains the translation-invariant caching that is the decisive optimisation
for a regular grid. Sections~\ref{sec:complexity}--\ref{sec:params} analyse
complexity and accuracy and justify every parameter. Sections~\ref{sec:valid}--\ref{sec:outlook}
report validation and discuss limitations. Throughout, claims of fact about the
implementation are tied to specific source locations.

\paragraph{Conventions.}
A grid element is identified by integer indices $(i_x,i_y)$, $0\le i_x,i_y<N_s$,
and stored at the global flat index $k=i_y N_s + i_x$ (row-major, $x$ fastest;
see \code{boussinesq\_kernel.hpp:30--34} and \code{h2\_operator.cpp:175}). The
element side is $h=L/N_s$ where $L$ is the domain size. We write $S_{ij}$ for the
influence of source element $j$ on target element $i$, $\xx$ for a target point,
$\yy$ for a source point, and $g$ for the (scaled) point kernel. Chebyshev order
is $q$; the per-box interpolation rank is $r=q^2$.

\section{The physical problem and its Green's function}
\label{sec:physics}

\subsection{Boussinesq's point solution and the displacement convolution}

Consider a linear-elastic, isotropic half-space occupying $z\le 0$ with the free
surface $z=0$. The classical Boussinesq solution gives the normal surface
displacement produced by a concentrated normal force $P$ applied at the origin of
the surface~\cite{Johnson1985}:
\begin{equation}
  u_z(r) = \frac{P}{\pi \Es\, r}, \qquad r=\sqrt{x^2+y^2},
  \label{eq:boussinesq-point}
\end{equation}
where the reduced (plane-strain composite) modulus is
\begin{equation}
  \frac{1}{\Es} = \frac{1-\nu_1^2}{E_1} + \frac{1-\nu_2^2}{E_2}
  \label{eq:estar}
\end{equation}
for two bodies with Young's moduli $E_{1,2}$ and Poisson ratios $\nu_{1,2}$; for a
rigid indenter on a single elastic body the second term drops. By linear
superposition, a distributed pressure $p(\yy)$ acting on the surface produces the
displacement field
\begin{equation}
  u_z(\xx) = \int_{\Omega} G(\xx-\yy)\, p(\yy)\, \mathrm{d}A_{\yy},
  \qquad
  G(\mathbf{r}) = \frac{1}{\pi \Es\, \abs{\mathbf{r}}}.
  \label{eq:convolution}
\end{equation}
Equation~\eqref{eq:convolution} is a two-dimensional convolution with the
half-space Green's function $G$; it is the surface-restriction of the
three-dimensional Boussinesq solution and is the only ingredient of the elastic
model that the solver needs.

\subsection{Why this kernel, and its analytical properties}

Three features of $G$ shape every algorithmic choice that follows.

\begin{enumerate}
\item \textbf{Singular but integrable.} $G\sim 1/\abs{\mathbf r}$ has a
non-integrable pointwise value at the origin but an integrable singularity in
2-D, so the convolution \eqref{eq:convolution} and, in particular, the
self-interaction of a finite element are finite. This is what lets a
piecewise-constant discretisation use an \emph{exact} closed-form self-term
rather than a regularisation.
\item \textbf{Smooth and analytic away from the origin.} For
$\abs{\mathbf r}>0$, $G$ is real-analytic. Between two well-separated regions the
kernel $G(\xx-\yy)$ is a smooth function of $(\xx,\yy)$, hence amenable to
polynomial interpolation with spectral accuracy
(Section~\ref{sec:degenerate}). This is the analytic fact that makes the FMM
work.
\item \textbf{Translation invariant.} $G$ depends only on the difference
$\xx-\yy$. On a uniform grid this collapses the entire $N\times N$ matrix to an
$N_s\times N_s$ table and makes the FMM transfer operators reusable across the
tree (Section~\ref{sec:cache}). This is the special structure that turns a
generic $\Order(N)$ FMM into one with a very small constant.
\end{enumerate}

The spectral view is also instructive. On the (bi-)infinite plane the convolution
operator $G\!*\!\cdot$ is diagonalised by the Fourier transform with symbol
\begin{equation}
  \widehat{G}(\mathbf q) = \frac{2}{\Es\,\abs{\mathbf q}},
  \label{eq:symbol}
\end{equation}
i.e.\ the operator behaves like $\abs{\mathbf q}^{-1}$: it is a smoothing
(order $-1$) pseudodifferential operator. The $1/\abs{\mathbf q}$ decay of the
symbol both explains why FFT methods are natural (Section~\ref{sec:cost},
\ref{sec:valid}) and why the off-diagonal blocks of $S$ are numerically low
rank (Section~\ref{sec:degenerate}).

\section{Discretisation: piecewise-constant pressure and the Love cell integral}
\label{sec:discretisation}

\subsection{Collocation with piecewise-constant pressure}

We tile the contact domain with a uniform $N_s\times N_s$ grid of square
elements of side $h=L/N_s$ and assume the pressure is constant over each element,
$p(\yy)=p_j$ for $\yy$ in element $j$. Collocating the displacement
\eqref{eq:convolution} at the centre $\xx_i$ of each target element $i$ gives the
linear system $u_i=\sum_j S_{ij} p_j$ with
\begin{equation}
  S_{ij} = \int_{\text{element }j} G(\xx_i-\yy)\,\mathrm{d}A_{\yy}.
  \label{eq:Sij-integral}
\end{equation}
The entry $S_{ij}$ is therefore the displacement at the centre of element $i$ due
to \emph{unit uniform pressure} on element $j$. The integral
\eqref{eq:Sij-integral} is over a square and can be evaluated in closed form.

\subsection{The Love (1929) closed-form integral}

Let a square source patch be centred at the origin and occupy
$[-a,a]\times[-b,b]$, carrying unit uniform pressure. The surface normal
displacement at a point $(x,y)$ is (Love~\cite{Love1929};
Johnson~\cite[eq.~3.25]{Johnson1985})
\begin{equation}
\begin{aligned}
  \pi\Es\, u_z(x,y;a,b) \;=\;
  & (x+a)\,\ln\!\frac{(y+b)+R_{++}}{(y-b)+R_{+-}}
  + (y+b)\,\ln\!\frac{(x+a)+R_{++}}{(x-a)+R_{-+}} \\[2pt]
  &+ (x-a)\,\ln\!\frac{(y-b)+R_{--}}{(y+b)+R_{-+}}
  + (y-b)\,\ln\!\frac{(x-a)+R_{--}}{(x+a)+R_{+-}},
\end{aligned}
\label{eq:love}
\end{equation}
where the four corner distances are
\begin{equation}
  R_{\pm\pm} = \sqrt{(x\pm a)^2 + (y\pm b)^2}.
  \label{eq:corner}
\end{equation}
This is implemented verbatim as \code{love\_uz(x,y,a,b)} in
\code{boussinesq\_kernel.cpp:5--16}; the function returns the right-hand side of
\eqref{eq:love} \emph{without} the $1/(\pi\Es)$ factor (the factor is applied by
the callers, see below). The four logarithm arguments and the four
\code{std::hypot} distances $R_{\pm\pm}$ in the code are exactly
\eqref{eq:love}--\eqref{eq:corner}.

\begin{remark}[Derivation in one line]
Equation~\eqref{eq:love} is obtained by integrating the point solution
\eqref{eq:boussinesq-point} over the rectangle,
$\int_{-a}^{a}\!\int_{-b}^{b} \big[(x-\xi)^2+(y-\eta)^2\big]^{-1/2}
\mathrm{d}\eta\,\mathrm{d}\xi$, using
$\int (s^2+c^2)^{-1/2}\mathrm{d}s = \operatorname{arcsinh}(s/c)
=\ln\!\big(s+\sqrt{s^2+c^2}\big)+\text{const}$ for each variable in turn. The
boundary terms evaluated at the four corners assemble into the four
$(\cdot)\ln(\cdot)$ groups of \eqref{eq:love}.
\end{remark}

\subsection{The self-term}

For Hcontact's square elements $a=b=h/2$. Setting $x=y=0$ in \eqref{eq:love}
gives the diagonal entry $S_{ii}$. All four corner distances equal
$R_{\pm\pm}=\tfrac{h}{2}\sqrt2$, and each of the four logarithm arguments reduces
to $(1+\sqrt2)/(\sqrt2-1)=(1+\sqrt2)^2$, so each term contributes
$\tfrac{h}{2}\cdot 2\ln(1+\sqrt2)$:
\begin{equation}
  S_{ii} = \frac{4h\,\ln(1+\sqrt2)}{\pi \Es}.
  \label{eq:selfterm}
\end{equation}
This closed-form self-term (no singular quadrature) is one practical benefit of
choosing the integrated Love kernel over the bare point kernel; it appears in
the table at offset $(0,0)$, \code{boussinesq\_kernel.cpp:23--27}.

\subsection{Translation invariance and the $N_s\times N_s$ lookup table}
\label{sec:table}

Because the grid is regular and $G$ depends only on $\xx-\yy$, the integral
\eqref{eq:Sij-integral} depends only on the index offset
$(\Delta_x,\Delta_y)=(\abs{i_x-j_x},\abs{i_y-j_y})$:
\begin{equation}
  S_{ij} = T[\Delta_y,\Delta_x],
  \qquad
  T[d_y,d_x] = \frac{1}{\pi\Es}\,\code{love\_uz}\!\big(d_x h,\, d_y h,\, \tfrac{h}{2},\tfrac{h}{2}\big).
  \label{eq:table}
\end{equation}
The constructor \code{BoussinesqKernel} builds this $N_s\times N_s$ table once
with $\Order(N)$ kernel evaluations (\code{boussinesq\_kernel.cpp:18--27}); the
sign symmetry $S_{ij}=S_{ji}$ and reflection symmetry justify storing only
non-negative offsets. Lookups are $\Order(1)$: \code{entry(i,j)} reduces the flat
indices to $(\Delta_x,\Delta_y)$ (\code{:30--34}), and \code{entry\_offset(dix,diy)}
serves a \emph{signed} offset, returning $0$ when the offset leaves the grid span
(\code{:38--42}); the latter is what the $\Htwo$ near-field stencil cache calls
(Section~\ref{sec:near}). The matrix $S$ is symmetric positive definite, which is
why the contact QP below is convex.

\subsection{The contact problem and the role of the operator}

The discrete contact problem is the bound- and equality-constrained quadratic
program
\begin{equation}
  \min_{\pp}\ \tfrac12\,\pp^\top S\,\pp + \pp^\top \mathbf g_0
  \quad\text{s.t.}\quad \pp\ge 0,\ \ \operatorname{mean}(\pp)=\bar p,
  \label{eq:qp}
\end{equation}
solved by the Polonsky--Keer projected conjugate gradient~\cite{PolKeer1999}
(documented in \code{doc/theory/pcg.tex}). \emph{Each PCG iteration needs two
applications of $S$} and never an individual entry, a factorisation, or the
stored matrix. The operator is therefore required \emph{only} as a
matrix--vector product $\mathbf u = S\pp$. This single observation selects the
class of algorithm: anything that delivers a fast, low-memory matvec is
admissible, and we are free to trade an exactly-formed $S$ for a controlled
approximation of the apply.

\section{The cost problem and the landscape of fast methods}
\label{sec:cost}

\subsection{Why $N$ gets large}

For self-affine rough surfaces the contact pressure has structure down to the
finest resolved wavelength, so converged results demand fine grids. With
$N=N_s^2$, the dense matrix needs $8N^2$ bytes: $128$\,MiB at $N_s=64$,
$2$\,GiB at $N_s=128$, $32$\,GiB at $N_s=256$, $512$\,GiB at $N_s=512$, and
$\approx 8.8$\,TB at $N_s=1024$. A dense apply is $\Order(N^2)$ flops per
iteration. Both are hopeless beyond a few hundred elements per side. The same
table also shows why even storing the dense matrix once for reuse is not an
option: it is the storage, not just the flops, that fails.

\subsection{Four representations}

Table~\ref{tab:landscape} summarises the candidates.

\begin{table}[h]
\centering
\begin{tabular}{lll}
\toprule
representation & memory & matvec \\
\midrule
dense matrix                       & $\Order(N^2)$       & $\Order(N^2)$ \\
classical $\Hmat$-matrix (ACA)     & $\Order(N\log N)$   & $\Order(N\log N)$ \\
FFT convolution (zero-padded)      & $\Order(N)$         & $\Order(N\log N)$ \\
$\Htwo$/FMM (this document)        & $\Order(N)$         & $\Order(N)$ \\
\bottomrule
\end{tabular}
\caption{Asymptotic cost of the four representations of the apply $\mathbf u=S\pp$.}
\label{tab:landscape}
\end{table}

\paragraph{Dense.} Exact, trivial to implement, and used in Hcontact only for
small-$N$ tests (\code{assemble\_dense}, \code{boussinesq\_kernel.cpp:29--36}). It
is the ground truth against which the $\Htwo$ apply is verified
(Section~\ref{sec:accuracy}).

\paragraph{Classical $\Hmat$-matrix with ACA.} A block-cluster tree partitions
$S$ into admissible (well-separated) and inadmissible blocks; each admissible
block $A_{ts}$ is compressed independently as a low-rank product
$A_{ts}\approx U_{ts}V_{ts}^\top$ by adaptive cross
approximation~\cite{Bebendorf2008,Hackbusch2015}. This works and was Hcontact's
previous backend, but for a \emph{translation-invariant} kernel it stores
\emph{redundant bases}: many admissible blocks at one level share the same
relative geometry and are therefore the \emph{same} matrix, yet each is
approximated and stored separately. Memory is $\Order(N\log N)$ with a large
constant (e.g.\ $6.2$\,GiB at $N_s=512$ in Hcontact). The $\Htwo$/FMM design
exists precisely to remove this redundancy.

\paragraph{FFT convolution.} Because $S$ is a (block-)Toeplitz convolution on a
regular grid, a zero-padded FFT applies it in $\Order(N\log N)$ with $\Order(N)$
memory; this is the discrete-convolution-FFT (DC-FFT) method of Stanley and
Kato~\cite{StanleyKato1997} and Liu, Wang and Liu~\cite{LiuWangLiu2000}, and is
what the reference solver Tamaas implements. It is excellent for a fully active
rectangular grid but is intrinsically \emph{global}: it cannot cheaply skip work
outside the contact patch, and it ties the algorithm to the regular grid and
periodic/zero-padded boundary handling. We compare against it quantitatively in
Section~\ref{sec:valid}.

\paragraph{$\Htwo$/FMM.} A nested (level-consistent) low-rank representation in
which each cluster carries a \emph{single} interpolation basis shared by all its
blocks, and inter-cluster coupling is a small matrix~\cite{FongDarve2009,Hackbusch2015,BoerHugesAnda2006}.
This is the only representation with $\Order(N)$ \emph{both} in memory and in
matvec, and -- crucially for this kernel -- its building blocks become reusable
caches under translation invariance. It is the subject of the rest of this
document.

\subsection{From $\Hmat$ to $\Htwo$: the nestedness idea}

In a classical $\Hmat$-matrix the rank-$k$ factor $U_{ts}$ of a block has its own
column space for every pair $(t,s)$. The $\Htwo$ idea is to demand a
\emph{nested} cluster basis: a single matrix $V_t$ per cluster $t$ (independent
of the partner $s$) such that every admissible block factors as
\begin{equation}
  A_{ts} \approx V_t\, K_{ts}\, V_s^\top,
  \label{eq:h2-factor}
\end{equation}
with a small coupling matrix $K_{ts}$, and such that a child basis is obtained
from its parent's basis by a fixed transfer matrix (the M2M operator). Two
consequences follow immediately: (i) bases are shared across all blocks of a
cluster, eliminating the ACA redundancy; and (ii) the apply can be organised as
an upward pass (form moments $V_s^\top \pp_s$ once per cluster), a coupling pass
(apply $K_{ts}$), and a downward pass (expand by $V_t$), which is the FMM. The
bbFMM realises the bases $V_t$ as polynomial-interpolation operators, so that
\eqref{eq:h2-factor} is exactly the degenerate-kernel approximation of
Section~\ref{sec:degenerate}, and $K_{ts}$ is simply the kernel sampled at
interpolation nodes.

\section{Hierarchical geometry: the uniform quadtree}
\label{sec:tree}

\subsection{Levels, boxes, and index ranges}

The grid is partitioned by a \emph{balanced, complete} quadtree
(\code{UniformQuadTree}, \code{uniform\_quadtree.\{hpp,cpp\}}). Level $0$ is the
root box covering the whole grid; each box at level $l$ has four children at
level $l+1$ obtained by bisecting in $x$ and $y$. With $N_s$ and the leaf side
$\ell$ (\code{leaf\_side}) both powers of two and $\ell\mid N_s$, the number of
levels is
\begin{equation}
  L_{\max}+1 = \log_2\!\big(N_s/\ell\big) + 1
  \label{eq:nlevels}
\end{equation}
(\code{uniform\_quadtree.cpp:21}), with the leaf level $L_{\max}=$
\code{leaf\_level()}. A box at level $l$ has side $s = N_s\gg l = N_s/2^{l}$
elements and box-coordinates $(b_x,b_y)\in[0,2^l)^2$; it covers the element range
\begin{equation}
  [\,b_x s,\ (b_x{+}1)s\,)\times[\,b_y s,\ (b_y{+}1)s\,)
\end{equation}
(\code{:36--44}). Crucially, a box stores only this \emph{index range}
$(i_{x0},i_{y0},\text{side})$ together with parent/child ids -- never an explicit
list of the elements it contains. The global flat index of element $(i_x,i_y)$ is
recovered arithmetically as $i_y N_s + i_x$, so a box is a handful of integers
(\code{H2Box}, \code{uniform\_quadtree.hpp:11--18}). The whole tree is
$\Order(N/\ell^2)$ tiny structs: $\sum_{l=0}^{L_{\max}} 4^l = (4^{L_{\max}+1}-1)/3
\approx \tfrac{4}{3}(N/\ell^2)$ boxes, laid out by level via the prefix sum
\code{level\_begin\_} (\code{:23--30}), so that \code{box\_id(level,bx,by)} is
$\Order(1)$ (\code{uniform\_quadtree.hpp:35--37}).

\begin{remark}[Index ranges vs index lists]
A generic FMM stores, per leaf, the list of particles it owns. Here the regular
grid lets us replace lists by ranges; this is what keeps tree memory negligible
and lets every leaf gather/scatter its $\ell^2$ DOFs by two nested loops over
$(l_x,l_y)$ (\code{h2\_operator.cpp:173--175, 239--247}).
\end{remark}

\subsection{Geometric extent and the $[-1,1]^2$ normalisation}
\label{sec:extent}

For interpolation, each box is mapped affinely to the reference square
$[-1,1]^2$. Hcontact uses the box's \emph{geometric extent} (its full side
$s\cdot h$ in physical units), not the bounding box of its element centres. The
$k$-th element centre of a side-$s$ box maps to
\begin{equation}
  \hat c_k = \frac{k - s/2 + \tfrac12}{s/2}, \qquad k=0,\dots,s-1,
  \label{eq:centers-norm}
\end{equation}
implemented as \code{centers\_norm(side)} in \code{h2\_operator.cpp:15--22}. The
first and last centres map to $\pm(1-1/s)$, i.e.\ they sit just inside $\pm1$ by
half an element, consistent with cell-centred data filling the full box.

\begin{proposition}[Exact half-scaling]
\label{prop:halfscale}
With the geometric-extent normalisation, the reference-coordinate map of a child
box is related to its parent's by an exact factor $1/2$ plus a fixed shift
$\pm\tfrac12$, independent of level.
\end{proposition}
\begin{proof}
A parent of side $s$ has half-extent $w=\tfrac{s h}{2}$ and centre $c$; a child
has side $s/2$, half-extent $w/2$, and centre $c\pm w/2$ (the $\pm$ being the
quadrant). A point at child-reference coordinate $\xi$ is at the physical point
$(c\pm w/2)+(w/2)\xi$, which in the parent reference frame is
$\big[(c\pm w/2)+(w/2)\xi - c\big]/w = \pm\tfrac12 + \tfrac12\xi$. The map
$\xi\mapsto \pm\tfrac12+\tfrac12\xi$ has no dependence on $s$ or $l$.
\end{proof}

Proposition~\ref{prop:halfscale} is the geometric reason the M2M/L2L transfer
operators collapse to four level-independent matrices
(Section~\ref{sec:m2m}); it is precisely the shift used in
\code{h2\_operator.cpp:47--53}, where \code{signx}, \code{signy}
$=\pm0.5$ and the child node is placed at $\text{sign}+0.5\,\xi$.

\subsection{Near/far split, admissibility, and interaction lists}

Two boxes at the same level are \emph{near neighbours} when their Chebyshev
(max-norm) box-distance is $\le r=$\code{near\_radius}: this is the
$(2r+1)\times(2r+1)$ block centred on the box, clipped to the grid
(\code{neighbors(box,r)}, \code{uniform\_quadtree.cpp:58--69}; the default $r=1$
gives the $3\times3$ block). The near field is computed with the \emph{exact}
kernel (Section~\ref{sec:near}). Everything else is handled approximately via the
FMM \emph{interaction list}.

\begin{definition}[Interaction list]
\label{def:ilist}
For a non-root box $t$ with parent $P(t)$, the interaction list is
\begin{equation}
  I(t) = \big\{\, s : s \text{ is a child of a near neighbour of } P(t),\ \
                    s \notin \mathcal N(t) \,\big\},
  \label{eq:ilist}
\end{equation}
where $\mathcal N(t)$ is the set of near neighbours of $t$. The root has
$I=\varnothing$.
\end{definition}

This is exactly \code{interaction\_list(box,near\_radius)} in
\code{uniform\_quadtree.cpp:71--86}: it loops over \code{neighbors(parent,r)},
collects their children, and removes those in \code{neighbors(box,r)}. Members of
$I(t)$ are \emph{admissible}: they are well-separated from $t$ (separated by at
least one box width at their level) yet their parents are near, so the pair was
\emph{not} resolved one level up. This is the standard FMM admissibility, and it
matches the geometric criterion
$\min(\operatorname{diam} t,\operatorname{diam} s)\le \eta\,\operatorname{dist}(t,s)$
of $\Hmat$-theory with $\eta=\Order(1)$ (here $\eta\approx 2$ for $r=1$); see the
discussion in Section~\ref{sec:params-eta}.

\subsection{Exact tiling: every pair once}
\label{sec:tiling}

We now prove the central correctness property of the near/far split: the near
field and the per-level interaction lists together account for every ordered
pair of leaves exactly once.

For a fixed leaf pair $(\tau,\sigma)$ let $t_l,s_l$ denote their ancestor boxes
at level $l$ (so $t_{L_{\max}}=\tau$, $s_{L_{\max}}=\sigma$), and let
$\delta_l=\max(\abs{b_x(t_l)-b_x(s_l)},\abs{b_y(t_l)-b_y(s_l)})$ be their box
distance. Call the pair \emph{adjacent} at level $l$ if $\delta_l\le r$.

\begin{lemma}[Monotone separation]
\label{lem:monotone}
If $t_l,s_l$ are non-adjacent at level $l$ (i.e.\ $\delta_l\ge r+1$), then all
their descendants are non-adjacent: $\delta_{l'}\ge r+1$ for every $l'>l$.
\end{lemma}
\begin{proof}
Child box-coordinates satisfy $b_x(t_{l+1})=2\,b_x(t_l)+c_x$ with
$c_x\in\{0,1\}$, and similarly for $s$ and for $y$. Hence
$\abs{b_x(t_{l+1})-b_x(s_{l+1})} = \abs{2(b_x(t_l)-b_x(s_l))+(c_x^t-c_x^s)}
\ge 2\abs{b_x(t_l)-b_x(s_l)} - 1$, because $c_x^t-c_x^s\in\{-1,0,1\}$. Taking the
$\max$ over the two axes, $\delta_{l+1}\ge 2\delta_l-1$. If $\delta_l\ge r+1$
then $\delta_{l+1}\ge 2(r+1)-1 = 2r+1\ge r+1$. Induction on $l'$ gives the claim.
\end{proof}

\begin{theorem}[Exact-once coverage]
\label{thm:coverage}
Every ordered pair of leaves $(\tau,\sigma)$ is accounted for exactly once,
either by the near field at the leaf level, or by the M2L step at one unique
level.
\end{theorem}
\begin{proof}
At level $0$ both ancestors are the root, so $\delta_0=0\le r$: the pair is
adjacent at the root. By Lemma~\ref{lem:monotone}, the sequence of truth values
``adjacent at level $l$'' is non-increasing in $l$ (once false it stays false).
Two cases.

\emph{(i) Adjacent at the leaf level} ($\delta_{L_{\max}}\le r$). Then $\sigma$
is a near neighbour of $\tau$ and the pair is computed by the near field,
\code{h2\_operator.cpp:235--243}; it is never in any interaction list (members of
$I(\cdot)$ are explicitly excluded from being near neighbours, Def.~\ref{def:ilist}),
so it is counted exactly once.

\emph{(ii) Non-adjacent at the leaf level.} Let $l^\star\in\{1,\dots,L_{\max}\}$
be the unique level where the truth value flips: adjacent at $l^\star-1$
($\delta_{l^\star-1}\le r$) and non-adjacent at $l^\star$
($\delta_{l^\star}\ge r+1$); uniqueness is exactly the monotonicity. At level
$l^\star$, the parents $P(t_{l^\star})=t_{l^\star-1}$ and
$P(s_{l^\star})=s_{l^\star-1}$ are near neighbours, while $t_{l^\star}$ and
$s_{l^\star}$ are not. By Def.~\ref{def:ilist}, $s_{l^\star}\in I(t_{l^\star})$,
so the pair's contribution is carried by the M2L coupling at level $l^\star$
(the source leaf's value flows up via M2M into $s_{l^\star}$, across via
$K_{t_{l^\star}s_{l^\star}}$, and down via L2L into $\tau$). At every level
$l<l^\star$ the ancestors are adjacent, so $s_l\notin I(t_l)$; at every level
$l>l^\star$ the ancestors are non-adjacent \emph{but their parents are also
non-adjacent} (by Lemma~\ref{lem:monotone} applied to level $l-1\ge l^\star$),
so $s_l\notin I(t_l)$ there as well. Hence the pair appears in exactly one
interaction list, at level $l^\star$.
\end{proof}

Theorem~\ref{thm:coverage} guarantees that the FMM apply computes
$\mathbf u=S\pp$ with no double counting and no omission; the only error is the
interpolation error on the far blocks (Section~\ref{sec:accuracy}), since the
near blocks use the exact kernel.

\section{Degenerate-kernel and Chebyshev interpolation theory}
\label{sec:degenerate}

\subsection{Why a smooth far block is numerically low rank}

Let $t$ and $s$ be well-separated boxes and consider the kernel restricted to
them, $g(\xx,\yy)$ for $\xx\in t$, $\yy\in s$. A \emph{degenerate} (separable)
approximation of rank $r$ is
\begin{equation}
  g(\xx,\yy)\approx \sum_{m=1}^{r} \phi_m(\xx)\,\psi_m(\yy).
  \label{eq:degenerate}
\end{equation}
If such an approximation holds with error $\varepsilon$, then the corresponding
matrix block $A_{ts}$, with entries $g(\xx_i,\yy_j)$, is within $\varepsilon$ (in
the entrywise sense, hence up to grid factors in norm) of a rank-$r$ matrix:
$A_{ts}\approx \Phi\,\Psi^\top$ with $\Phi_{im}=\phi_m(\xx_i)$,
$\Psi_{jm}=\psi_m(\yy_j)$. Thus \emph{smoothness of $g$ on a separated pair
implies numerical low rank of the block}. For the analytic half-space kernel the
relevant approximation is polynomial interpolation, and the rank is the number
of interpolation nodes.

\subsection{Chebyshev nodes and the interpolation operator}
\label{sec:cheb}

The bbFMM realises \eqref{eq:degenerate} by interpolating $g$ in each argument at
Chebyshev points. Hcontact uses the Chebyshev points of the first kind on
$[-1,1]$,
\begin{equation}
  \xi_m = \cos\!\Big(\frac{(2m+1)\pi}{2q}\Big), \qquad m=0,\dots,q-1,
  \label{eq:cheb-nodes}
\end{equation}
the roots of the degree-$q$ Chebyshev polynomial $T_q$ (\code{ChebBasis},
\code{cheb\_basis.cpp:7--10}). For a function $f$ on $[-1,1]$ the degree-$(q{-}1)$
interpolant through these nodes can be written with the closed-form Lagrange-type
weights
\begin{equation}
  f(t) \approx \sum_{m=0}^{q-1} S_q(t,\xi_m)\, f(\xi_m),
  \qquad
  S_q(t,\xi_m) = \frac{1}{q} + \frac{2}{q}\sum_{n=1}^{q-1} T_n(\xi_m)\,T_n(t),
  \label{eq:Sq}
\end{equation}
where $T_n(\cos\theta)=\cos(n\theta)$ are the Chebyshev polynomials. This is the
weight returned by \code{ChebBasis::weights(t)} (\code{cheb\_basis.cpp:13--36}),
which builds the vector $w_m=S_q(t,\xi_m)$ by the three-term recurrence for
$T_n(\xi_m)$ and $T_n(t)$. Equation~\eqref{eq:Sq} is the discrete-cosine form of
the interpolation operator and is \emph{exact for all polynomials of degree
$<q$} (it reproduces each $T_n$, $n<q$, by orthogonality of the Chebyshev
polynomials at the first-kind nodes).

\begin{remark}[Interpolation operator, not the kernel]
$S_q$ is a property of the node set, not of $g$. It supplies the smooth
``shape functions'' $\phi_m(\xx)=S_q(\hat x,\xi_m)$ in \eqref{eq:degenerate}.
This is precisely what makes the method \emph{black-box}: the basis is
kernel-independent; only the coupling samples $g(\xi_a^t,\xi_b^s)$ depend on $g$.
\end{remark}

\subsection{Near-minimax property and spectral convergence}

Chebyshev interpolation is near-optimal among polynomial interpolants. Its
Lebesgue constant grows only logarithmically,
$\Lambda_q = \tfrac{2}{\pi}\ln q + \Order(1)$, so the interpolation error is
within a factor $\Lambda_q$ of the best degree-$(q{-}1)$ polynomial
approximation~\cite{Trefethen2013}. For analytic functions the best
approximation -- hence the interpolation error -- decays \emph{geometrically} in
$q$.

\begin{theorem}[Spectral convergence for analytic kernels]
\label{thm:cheb-conv}
Let $f$ be analytic on $[-1,1]$ and analytically continuable to the Bernstein
ellipse $E_\rho$ (foci $\pm1$, semi-axis sum $\rho>1$), with $\abs{f}\le M$ on
$E_\rho$. Then the degree-$(q{-}1)$ Chebyshev interpolant $p_{q-1}$ satisfies
\begin{equation}
  \norm{f-p_{q-1}}_{\infty,[-1,1]} \;\le\; \frac{4M}{\rho-1}\,\rho^{-(q-1)}.
  \label{eq:cheb-bound}
\end{equation}
\end{theorem}
This is the standard Chebyshev-interpolation estimate
(Trefethen~\cite{Trefethen2013}, Thm.~8.2). For two boxes separated by at least
their own width, the kernel $g(\xx,\yy)=1/(\pi\Es\abs{\xx-\yy})$ extends
analytically in each variable to an ellipse with $\rho>1$ bounded away from $1$
by the separation ratio; the larger the separation, the larger $\rho$, and the
faster the decay. The practical consequence is the geometric error reduction in
$q$ measured for Hcontact: relative $L_2$ matvec error
$1.3\times10^{-4}$ at $q=4$ falling to $3.2\times10^{-6}$ at $q=6$
(Section~\ref{sec:accuracy}), i.e.\ roughly a factor $40$ per two extra orders,
the signature of a geometric rate.

\subsection{Tensor product in 2-D and the rank $r=q^2$}

In two dimensions the interpolation operator is the tensor product of the 1-D
operator on each axis. For a target box, a point $\xx=(x,y)$ with reference
coordinates $(\hat x,\hat y)$ has the 2-D weights
\begin{equation}
  S_q^{(2)}\big(\xx,\boldsymbol\xi_A\big)
  = S_q(\hat x,\xi_{A_x})\,S_q(\hat y,\xi_{A_y}),
  \qquad A=(A_x,A_y),\ \ A_x,A_y\in\{0,\dots,q-1\},
  \label{eq:tensor}
\end{equation}
so the node set is the $q\times q$ Chebyshev grid $\boldsymbol\xi_A=(\xi_{A_x},
\xi_{A_y})$, and the per-box rank is
\begin{equation}
  r = q^2
\end{equation}
(\code{q2\_ = q*q}, \code{h2\_operator.cpp:11}). The tensor structure is the key
to efficiency: rather than forming an $r\times\ell^2$ matrix and applying it as a
Kronecker product, each operator is applied as \emph{two sequential 1-D passes},
turning $\Order(r\,\ell^2)$ work into $\Order(q\,\ell^2 + q^2\ell)$ -- see the
$W^\top X W$ form in Section~\ref{sec:p2m}.

\section{The black-box FMM: derivation of the five transfer operators}
\label{sec:bbfmm}

We now derive P2M, M2M, M2L, L2L, L2P from the interpolation
\eqref{eq:tensor}--\eqref{eq:Sq}, and exhibit the small-matrix forms exactly as
they appear in \code{h2\_operator.cpp}. Throughout, capital indices $A,B$ run over
the $q^2$ Chebyshev nodes of a box; lowercase $l_x,l_y$ over the $\ell^2$ leaf
DOFs.

\subsection{The far interaction in interpolated form}

Fix an admissible target/source pair $(t,s)$ at a common level, with reference
maps $\xx=c_t+w_t\boldsymbol\xi$ and $\yy=c_s+w_s\boldsymbol\xi$ ($w$ the
half-extent, $c$ the centre). Interpolating $g$ in \emph{both} arguments on the
node grid gives the degenerate approximation
\begin{equation}
  g(\xx,\yy)\;\approx\;
  \sum_{A}\sum_{B} S_q^{(2)}(\xx,\boldsymbol\xi_A^{\,t})\,
  g\big(\boldsymbol\xi_A^{\,t},\boldsymbol\xi_B^{\,s}\big)\,
  S_q^{(2)}(\yy,\boldsymbol\xi_B^{\,s}),
  \label{eq:far-interp}
\end{equation}
where $\boldsymbol\xi_A^{\,t}=c_t+w_t\boldsymbol\xi_A$ are the physical node
positions of box $t$. Substituting into the far part of the apply,
$u(\xx_i)=\sum_{j\in s} g(\xx_i,\yy_j)\,p_j$, and regrouping:
\begin{equation}
  u(\xx_i)\ \approx\
  \underbrace{\sum_A S_q^{(2)}(\xx_i,\boldsymbol\xi_A^{\,t})}_{\text{L2P}}
  \;\underbrace{\sum_B g(\boldsymbol\xi_A^{\,t},\boldsymbol\xi_B^{\,s})}_{\text{M2L}}
  \;\underbrace{\sum_{j\in s} S_q^{(2)}(\yy_j,\boldsymbol\xi_B^{\,s})\,p_j}_{\text{P2M}}.
  \label{eq:regroup}
\end{equation}
The three braces are exactly the three leaf-level FMM operators; the upward and
downward tree passes (M2M, L2L) merely move the moment/local coefficients between
levels so that each $(t,s)$ pair is touched at its $l^\star$ (Theorem~\ref{thm:coverage}).

\subsection{P2M: pressure to moments (anterpolation)}
\label{sec:p2m}

The moment of source leaf $s$ at node $B$ is
$M_B^s = \sum_{j\in s} S_q^{(2)}(\yy_j,\boldsymbol\xi_B^{\,s})\,p_j$. Using the
tensor split \eqref{eq:tensor} and arranging the $\ell^2$ leaf values as a matrix
$X_s$ with $X_s[l_y,l_x]=p_{(l_x,l_y)}$,
\begin{equation}
  M^s = W^\top X_s\, W, \qquad W \equiv W_{\text{leaf}}\in\R^{\ell\times q},\quad
  W[l,a]=S_q(\hat c_l,\xi_a),
  \label{eq:p2m}
\end{equation}
where $\hat c_l$ are the normalised leaf centres \eqref{eq:centers-norm}. The
matrix $W_{\text{leaf}}$ is built once as
\code{cheb\_.weights\_at(centers\_norm(ls\_))} (\code{h2\_operator.cpp:43}), and
the apply is \code{Wleaf\_.transpose() * Xs * Wleaf\_} returning a $q\times q$
block (\code{:176}), then flattened to the length-$q^2$ vector
$M_t(a_x+q\,a_y)=M^{\text{mat}}(a_y,a_x)$ (\code{:177--180}). The name
\emph{anterpolation} reflects that this is the transpose of interpolation: it
``spreads'' the point values onto polynomial coefficients.

\subsection{M2M: child moments to parent (upward pass)}
\label{sec:m2m}

A parent's moments are obtained by re-interpolating its children's moments. The
$B$-th node of child $c$ sits, in the parent's reference frame, at
$\boldsymbol\xi^{\,(\text{par})}=\mathbf{sign}_c/2 + \tfrac12\boldsymbol\xi_B$
(Proposition~\ref{prop:halfscale}), where $\mathbf{sign}_c=(\pm\tfrac12,\pm\tfrac12)$
selects the quadrant. Re-interpolating onto the parent nodes,
\begin{equation}
  M^{\text{par}}_A \mathrel{+}= \sum_{B} R^{(c)}_{A,B}\, M^{\text{child}}_B,
  \qquad
  R^{(c)}_{A,B} = w_{A_x}\!\big(\tfrac{\text{sign}^x_c}{1}+\tfrac12\xi_{B_x}\big)\,
                  w_{A_y}\!\big(\tfrac{\text{sign}^y_c}{1}+\tfrac12\xi_{B_y}\big),
  \label{eq:m2m}
\end{equation}
where $w_{A}(\cdot)=S_q(\cdot,\xi_A)$. In the code (\code{h2\_operator.cpp:46--64})
this is built per child quadrant $c\in\{0,1,2,3\}$ with
\code{signx=(c\&1)?0.5:-0.5}, \code{signy=(c\&2)?0.5:-0.5}, the per-axis weight
vectors \code{Wx[a]=cheb\_.weights(signx+0.5*nodes[a])} (\code{:52--53}), and the
filled matrix \code{R(A, ax+q*ay)=Wx[ax](Ax)*Wy[ay](Ay)} with $A=A_x+q A_y$
(\code{:55--62}). The apply is \code{acc += R\_[c] * M[child]}
(\code{:191--195}), summed over the four children, sweeping from the finest
non-leaf level up to the root (\code{:184--197}).

By Proposition~\ref{prop:halfscale} the four matrices $R^{(c)}$ are
\emph{independent of level}: there are exactly four $q^2\times q^2$ M2M matrices,
reused at every level. This is the first half of the translation-invariant
collapse (Section~\ref{sec:cache}).

\subsection{M2L: source to target coupling (the kernel-dependent step)}
\label{sec:m2l}

The local expansion at a target box accumulates, over its interaction list, the
kernel sampled at node pairs:
\begin{equation}
  L^t_A \mathrel{+}= \sum_{s\in I(t)}\sum_B K^{ts}_{A,B}\, M^s_B,
  \qquad
  K^{ts}_{A,B} = g\big(\boldsymbol\xi_A^{\,t}-\boldsymbol\xi_B^{\,s}\big).
  \label{eq:m2l}
\end{equation}
Because $g$ is translation invariant, the coupling depends only on the node
\emph{difference}. With node coordinate $=$ centre $+\,b_h\,\xi$ where
$b_h=\tfrac12\,\text{side}\,h$ is the box half-extent, and the centre offset
between $t$ and $s$ equal to $(\Delta b_x,\Delta b_y)\cdot\text{side}\cdot h$, the
argument is
\begin{equation}
  \mathbf D = \big(\Delta b_x\,\text{side}\,h + b_h(\xi_{A_x}-\xi_{B_x}),\
                   \Delta b_y\,\text{side}\,h + b_h(\xi_{A_y}-\xi_{B_y})\big),
  \quad K^{ts}_{A,B}=g(\mathbf D).
  \label{eq:m2l-D}
\end{equation}
This is built verbatim in \code{h2\_operator.cpp:80--96}: \code{Dx = dbx*side*h +
bh*(nodes[Ax]-nodes[Bx])}, \code{Dy} analogous, with
\code{K(A, Bx+q*By) = far\_kernel(Dx,Dy)} and $b_h=$\code{0.5*side*h}. The far
point kernel is \code{far\_kernel(dx,dy) = scale * love\_uz(dx,dy,0.5h,0.5h)},
\code{scale} $=1/(\pi\Es)$ (\code{:24--26}) -- i.e.\ the \emph{same} Love cell
integral as the near field, evaluated at continuous node offsets
(Section~\ref{sec:lovefar}). The apply is grouped by target,
\code{acc += couplings\_[fi.coupling\_id] * M[fi.source\_box]}
(\code{:200--206}).

\begin{remark}[Why M2L is the only kernel-dependent operator]
P2M, M2M, L2L, L2P are assembled purely from the interpolation operator $S_q$ and
the box geometry; they would be identical for \emph{any} kernel. Only M2L
evaluates $g$. This is the defining structure of bbFMM: the kernel enters through
a single family of small matrices $K^{ts}$, which is what makes the method
black-box (works for any smooth $g$ via evaluations) and what makes the
translation-invariant caching so effective (Section~\ref{sec:cache}).
\end{remark}

\subsection{L2L: parent local to children (downward pass)}

L2L is the adjoint of M2M: the same transfer matrix moves local coefficients the
other way, parent to child,
\begin{equation}
  L^{\text{child}}_B \mathrel{+}= \sum_A R^{(c)}_{A,B}\, L^{\text{par}}_A
  \;=\; \big(R^{(c)\top} L^{\text{par}}\big)_B.
  \label{eq:l2l}
\end{equation}
In code, the child's quadrant is recovered from its parity
\code{c = (bx\&1) + 2*(by\&1)} and the update is
\code{L[b] += R\_[c].transpose() * L[parent]} (\code{h2\_operator.cpp:209--218}),
sweeping from level $1$ down to the leaves. The use of the \emph{transpose} of
the very same $R^{(c)}$ that drives M2M is what makes the whole operator
symmetric (Section~\ref{sec:complexity}, Remark~\ref{rem:symmetry}).

\subsection{L2P: target local to displacement}

Finally the local expansion is evaluated at the target leaf's DOFs. With the
local coefficients reshaped to a $q\times q$ matrix $L^{\text{mat}}$, the tensor
interpolation back to the $\ell\times\ell$ leaf grid is
\begin{equation}
  Y_t = W\, L^{\text{mat}}\, W^\top, \qquad
  u_{(l_x,l_y)} \mathrel{+}= Y_t[l_y,l_x],
  \label{eq:l2p}
\end{equation}
implemented as \code{Yt = Wleaf\_ * Lmat * Wleaf\_.transpose()}
(\code{h2\_operator.cpp:226--231}). Note that L2P uses $W$ where P2M used
$W^\top$: again the two leaf operators are mutual transposes.

\subsection{The near field}
\label{sec:near}

For each leaf $t$ and each near neighbour $s\in\mathcal N(t)$ the exact
contribution is a dense $\ell^2\times\ell^2$ block of $S$,
\begin{equation}
  A^{\text{near}}_{(l_x^t,l_y^t),(l_x^s,l_y^s)}
  = T\big[\,\abs{(l_y^t-l_y^s)-\Delta y\,\ell}\,,\ \abs{(l_x^t-l_x^s)-\Delta x\,\ell}\,\big],
  \label{eq:near}
\end{equation}
served from the lookup table $T$ via \code{kernel\_->entry\_offset(lxt - lxs -
dx*ls, lyt - lys - dy*ls)}, where $(\Delta x,\Delta y)=(b_x(s)-b_x(t),
b_y(s)-b_y(t))$ is the source-minus-target leaf offset
(\code{h2\_operator.cpp:113--130}). The apply gathers the source leaf's
$\ell^2$ values and does a dense matvec
\code{yloc += near\_stencils\_[ni.stencil\_id] * xs}
(\code{:235--243}), added to the L2P result (\code{:245--247}). Because the near
field uses the exact table entries, it contributes \emph{zero} approximation
error; all error is confined to the far interpolation.

\subsection{The six passes, assembled}

A matvec $\mathbf u = S\pp$ is, in order
(\code{h2\_operator.cpp:157--251}),
\[
  \text{P2M}\ \longrightarrow\ \text{M2M}\!\uparrow\ \longrightarrow\
  \text{M2L}\ \longrightarrow\ \text{L2L}\!\downarrow\ \longrightarrow\
  \text{L2P}\ \;+\;\ \text{near}.
\]
P2M fills the leaf moments $M$; M2M accumulates them up the tree; M2L produces
the local coefficients $L$ from far interactions; L2L pushes them down to the
leaves; L2P evaluates them; and the near pass adds the exact short-range part. The
moment/local buffers $M,L\in\R^{q^2}$ are allocated per box (\code{:160--164}).
All passes are OpenMP-parallel over boxes or leaves; the leaf passes
(P2M and L2P+near) write disjoint output ranges and so need no synchronisation
(\code{:168, 187, 200, 212, 221}).

\section{Translation-invariant caching: the decisive optimisation}
\label{sec:cache}

The bbFMM as derived is already $\Order(N)$ in work, but a \emph{generic} bbFMM
stores M2L matrices per interaction and M2M matrices per box, giving $\Order(N)$
or $\Order(N\log N)$ \emph{stored} small matrices with a sizeable constant. On a
uniform grid with a translation-invariant kernel, three collapses reduce this to
$\Order(\log N)+\Order(1)$ distinct matrices.

\begin{proposition}[M2M/L2L collapse]
\label{prop:m2m-cache}
With geometric-extent normalisation, there are exactly four distinct M2M matrices
$R^{(0)},\dots,R^{(3)}$ (one per child quadrant), independent of level, and the
four L2L matrices are their transposes.
\end{proposition}
\begin{proof}
By Proposition~\ref{prop:halfscale} the child-to-parent reference map is
$\xi\mapsto \mathbf{sign}_c/2+\tfrac12\xi$ for every level. Substituting into
\eqref{eq:m2m} shows $R^{(c)}$ depends only on $c$ and $q$, not on the level or
the box. There are four quadrants, hence four matrices. L2L uses $R^{(c)\top}$
by \eqref{eq:l2l}.
\end{proof}
In the code these four matrices are the array \code{R\_} of length $4$, built once
in \code{build()} (\code{h2\_operator.cpp:46--64}).

\begin{proposition}[M2L collapse]
\label{prop:m2l-cache}
The coupling $K^{ts}$ depends only on the triple $(\text{level},\Delta b_x,
\Delta b_y)$. The number of distinct couplings is $\Order(\log N)$.
\end{proposition}
\begin{proof}
From \eqref{eq:m2l-D}, $K^{ts}$ is a function of $b_h$ (set by the level via
$\text{side}=N_s\gg l$) and the integer offset $(\Delta b_x,\Delta b_y)$ only;
translation invariance removes any dependence on the absolute box position. The
interaction list \eqref{eq:ilist} has a \emph{bounded} number of distinct
offsets per level (for $r=1$ at most $(2\cdot3+1)^2-(2\cdot1+1)^2$ patterns,
i.e.\ a constant), and there are $L_{\max}+1=\Order(\log N)$ levels; hence
$\Order(\log N)$ distinct couplings.
\end{proof}
The cache is keyed by the packed integer \code{far\_key(level,dbx,dby)}
(\code{h2\_operator.cpp:29--33}); each new key creates one $q^2\times q^2$ matrix
and stores its id, reused by all interactions with the same key
(\code{:68--102}). The number of stored couplings is \code{n\_unique\_couplings}.

\begin{proposition}[Near-stencil collapse]
\label{prop:near-cache}
There are at most $(2r+1)^2$ distinct near stencils, indexed by the relative leaf
offset $(\Delta x,\Delta y)$; this is $\Order(1)$.
\end{proposition}
\begin{proof}
By \eqref{eq:near} the near block depends only on $(\Delta x,\Delta y)$, which
ranges over the $(2r+1)\times(2r+1)$ neighbourhood. Boundary leaves see a subset
of these offsets, so the count is at most $(2r+1)^2$.
\end{proof}
The cache is keyed by \code{near\_key(dx,dy)} (\code{h2\_operator.cpp:34--36}),
storing each $\ell^2\times\ell^2$ block once (\code{:109--138}); the count is
\code{n\_near\_stencils} (for $r=1$, nine interior stencils, the constant
$0.28$\,MiB cache reported in the concise note).

\begin{remark}[Contrast with classical $\Hmat$]
A classical $\Hmat$-matrix stores an independent $U_{ts},V_{ts}$ pair for every
admissible block, i.e.\ $\Order(N\log N)$ floating-point factors. The three
propositions replace this by $\Order(\log N)$ couplings $+\,\Order(1)$ stencils
$+\,4$ transfer matrices. This is the structural reason the $\Htwo$ operator
attains $\Order(N)$ \emph{memory} (the only $\Order(N)$ term being the
coefficient buffers $M,L$), with a small constant.
\end{remark}

\subsection{The Love far kernel, not $1/r$}
\label{sec:lovefar}

The far point kernel is taken to be the Love cell integral
$g(\mathbf d)=\code{love\_uz}(\mathbf d,h/2,h/2)/(\pi\Es)$
(\code{h2\_operator.cpp:24--26}), \emph{not} the bare point kernel
$1/(\pi\Es\abs{\mathbf d})$. The reason is consistency: the near field uses the
exact integrated entries $T$, so using the \emph{same} integrated kernel for the
far field makes the near and far representations of $S$ agree on the underlying
operator. Consequently the only far-field error is the interpolation error of
$g$ (Theorem~\ref{thm:cheb-conv}); there is no mismatch between a ``point'' far
model and an ``integrated'' near model. Since $g$ is just as smooth as $1/r$ away
from the origin (it differs by an $\Order(h^2)$ smoothing at the scale of one
cell, negligible at admissible separations), the interpolation rate is unchanged.

\section{Complexity and memory}
\label{sec:complexity}

\subsection{Operation counts}

Let $n_{\text{box}}\approx\tfrac43 N/\ell^2$ and $n_{\text{leaf}}=N/\ell^2$. Each
pass costs:
\begin{itemize}
\item \textbf{P2M / L2P}: two 1-D contractions $W^\top X W$ per leaf, costing
$\Order(q\ell^2 + q^2\ell)$; total $\Order\!\big((N/\ell^2)(q\ell^2+q^2\ell)\big)
= \Order(Nq + Nq^2/\ell)=\Order(N)$.
\item \textbf{M2M / L2L}: up to four $q^2\times q^2$ matvecs per box,
$\Order(q^4)$ each; total $\Order(n_{\text{box}}\,q^4)=\Order(N\,q^4/\ell^2)
=\Order(N)$.
\item \textbf{M2L}: $\Order(q^4)$ per far interaction; the number of far
interactions is $\Order(n_{\text{box}})$ because each interaction list has bounded
size; total $\Order(N\,q^4/\ell^2)=\Order(N)$.
\item \textbf{Near}: up to $(2r+1)^2$ dense $\ell^2\times\ell^2$ matvecs per leaf,
$\Order(\ell^4)$ each; total $\Order(n_{\text{leaf}}(2r+1)^2\ell^4)
=\Order\!\big(N(2r+1)^2\ell^2\big)=\Order(N)$.
\end{itemize}
Every pass is $\Order(N)$, so the matvec is $\Order(N)$. The near pass carries the
largest constant ($\propto (2r+1)^2\ell^2$), which is why $\ell$ and $r$ are kept
small (Section~\ref{sec:params}).

\subsection{Storage}

The exact byte counts are computed in \code{h2\_operator.cpp:151--154}:
\begin{equation}
\begin{aligned}
  \text{bytes}_{\text{coupling}} &= 8\,n_{\text{couplings}}\,q^4
     &&= \Order(q^4\log N),\\
  \text{bytes}_{\text{near}} &= 8\,n_{\text{stencils}}\,\ell^4
     &&= \Order(\ell^4) = \Order(1),\\
  \text{bytes}_{\text{buffers}} &= 8\cdot 2\,n_{\text{box}}\,q^2
     &&= \Order(N\,q^2/\ell^2) = \Order(N).
\end{aligned}
\label{eq:bytes}
\end{equation}
The coefficient buffers dominate asymptotically. Their cost per DOF is
\begin{equation}
  \frac{\text{bytes}_{\text{buffers}}}{N}
  = \frac{8\cdot2\cdot\tfrac43 (N/\ell^2)\,q^2}{N}
  = \frac{64}{3}\,\frac{q^2}{\ell^2}\ \text{bytes/DOF},
\end{equation}
which for the defaults $q=6,\ell=8$ is $\tfrac{64}{3}\cdot\tfrac{36}{64}=12$
bytes/DOF, matching the $\approx 13$ bytes/DOF asymptote quoted in the concise
note (the small excess is the $\Order(\log N)$ coupling cache). Total memory is
$\Order(N)$ with a tiny constant: e.g.\ $1.12$\,MiB at $N_s=64$, $5.26$\,MiB at
$N_s=512$, and $51$\,MiB at $N_s=2048$, versus $128$\,TiB for the dense matrix at
$N_s=2048$ (Table~\ref{tab:mem}). The near-stencil cache is a constant
$0.28$\,MiB.

\begin{table}[h]
\centering
\begin{tabular}{rrrrrr}
\toprule
$N_s$ & $N$ & coupling & buffers & total & dense $N^2$ \\
      &     & [MiB]    & [MiB]   & [MiB] & [MiB] \\
\midrule
64   & 4\,096      & 0.79 & 0.05 & 1.12  & 128 \\
128  & 16\,384     & 1.19 & 0.19 & 1.66  & 2\,048 \\
256  & 65\,536     & 1.58 & 0.75 & 2.61  & 32\,768 \\
512  & 262\,144    & 1.98 & 3.00 & 5.26  & 524\,288 \\
1024 & 1\,048\,576 & 2.37 & 12.0 & 14.65 & 8\,388\,608 \\
2048 & 4\,194\,304 & 2.77 & 48.0 & 51.05 & 134\,217\,728 \\
\bottomrule
\end{tabular}
\caption{$\Htwo$ operator memory ($q=6$, $\ell=8$, $r=1$) from
\code{bench\_h2\_memory.py}, reproduced from the concise note. The coupling cache
grows like $\log N$, the buffers like $N$; total memory crosses below the dense
cost by orders of magnitude.}
\label{tab:mem}
\end{table}

\begin{remark}[Symmetry of the operator]
\label{rem:symmetry}
Because L2P $=W(\cdot)W^\top$ is the transpose of P2M $=W^\top(\cdot)W$, L2L
$=R^{(c)\top}$ is the transpose of M2M $=R^{(c)}$, the near stencils are
symmetric ($T$ is even in each offset), and M2L is symmetric under
$t\leftrightarrow s$ (which negates $\mathbf D$, and $g$ is even), the assembled
$\Htwo$ apply is a \emph{symmetric} linear operator -- consistent with the SPD
matrix $S$ it approximates, and required by the conjugate-gradient solver.
\end{remark}

\paragraph{Allocation-free application (July 2026).}
The buffers of Table~\ref{tab:mem} -- the multipole and local coefficient
tables, $2\,q^2 n_{\text{box}}$ scalars -- are stored as two flat
$q^2\times n_{\text{box}}$ matrices owned by the operator and reused across
applications; the original implementation reallocated them as
$\sim\!2\,n_{\text{box}}$ individual small vectors on \emph{every} matvec,
which at $N_s=16384$ ($n_{\text{box}}\approx1.4\times10^6$) meant millions of
heap allocations per operator application. The per-box interaction lists are
flattened to CSR (offset array $+$ one contiguous entry array), the per-leaf
temporaries of P2M/L2P live in per-thread scratch inside the OpenMP regions,
and \code{matvec\_into(x, y)} writes into a caller-owned output vector, so a
steady-state matvec performs no heap allocation. Both precisions keep separate
persistent scratch (\code{double}/\code{float}), sized lazily on first use;
consequently a single \code{H2Operator} must not be applied from two threads
concurrently.

\subsection{Measured scaling}

Table~\ref{tab:cputime} (from the concise note, \code{bench\_h2\_cputime.py},
$q=6$, $20$-core node) confirms the linear scaling of build and matvec up to
$N_s=16384$ ($N\approx2.7\times10^8$). Each $4\times$ increase in $N$ multiplies
both build and matvec time by $\approx4$, the signature of $\Order(N)$. The
largest operator builds in $14$\,s, applies in $6$\,s, and fits in $14$\,GiB --
about $57$ bytes/DOF peak RSS (of which $\approx12$ are the intrinsic buffers, the
rest interaction-list indices and interpreter overhead).

\begin{table}[h]
\centering
\begin{tabular}{rrrrrrr}
\toprule
$N_s$ & $N$ & build [s] & matvec [s] & RSS [GiB] & solve [s] & iters \\
\midrule
128   & 16\,384       & 0.007 & 0.0005 & 0.04 & 0.21 & 90 \\
256   & 65\,536       & 0.011 & 0.0013 & 0.05 & 0.80 & 199 \\
512   & 262\,144      & 0.023 & 0.0052 & 0.08 & 3.64 & 241 \\
1024  & 1\,048\,576   & 0.064 & 0.019  & 0.19 & 48.0 & 906 \\
2048  & 4\,194\,304   & 0.245 & 0.098  & 0.26 & ---  & --- \\
4096  & 16\,777\,216  & 0.92  & 0.34   & 0.92 & ---  & --- \\
8192  & 67\,108\,864  & 3.65  & 1.39   & 3.57 & ---  & --- \\
16384 & 268\,435\,456 & 13.98 & 6.06   & 14.2 & ---  & --- \\
\bottomrule
\end{tabular}
\caption{$\Htwo$ CPU-time study ($q=6$, all cores) reproduced from the concise
note. Build and per-matvec scale $\approx\Order(N)$. Full rough-surface solves are
listed to $N_s=1024$; beyond that the projected-CG iteration count, not the
operator, becomes the limiting cost, so only operator metrics are reported.}
\label{tab:cputime}
\end{table}

\section{Accuracy analysis}
\label{sec:accuracy}

\subsection{Where the error lives}

By Theorem~\ref{thm:coverage} the near field is exact; by
Section~\ref{sec:lovefar} the far field uses the same integrated kernel as the
near field. Therefore the \emph{only} error in the apply is the tensor-Chebyshev
interpolation error of the far blocks, governed by Theorem~\ref{thm:cheb-conv}.
This error is uniform in $N$ (it depends on $q$ and the admissibility separation,
not on the grid size), so the relative matvec accuracy is essentially constant as
$N$ grows -- exactly what makes the method trustworthy at the largest grids,
where no dense ground truth is available.

\subsection{Measured interpolation error versus $q$}

Against the exact dense Love operator at $N_s=32$
(\code{tests/test\_h2.cpp}) the relative $L_2$ matvec error is
\begin{center}
\begin{tabular}{ccc}
\toprule
$q$ & rank $r=q^2$ & rel.\ $L_2$ matvec error \\
\midrule
4 & 16 & $1.3\times10^{-4}$ \\
6 & 36 & $3.2\times10^{-6}$ \\
\bottomrule
\end{tabular}
\end{center}
The roughly $40\times$ drop for two extra orders is the geometric rate predicted
by \eqref{eq:cheb-bound}. Plugged into the same PCG, the $\Htwo$ backend
reproduces the Hertz contact radius and peak pressure
($a_{\text{num}}/a_{\text{Hertz}}=1.016$, $p_{\max}/p_0=0.998$, from the
project's validated numbers) and matches the classical $\Hmat$-matrix contact
area.

\subsection{Effect of the near radius}

Increasing $r$ enlarges the exact near field and shrinks the interpolated far
field, so it can only \emph{improve} far accuracy (admissible pairs become more
separated, raising the Bernstein $\rho$). But near cost grows like $(2r+1)^2$
(Section~\ref{sec:complexity}). With the Love far kernel the default $r=1$ already
delivers $3\times10^{-6}$ at $q=6$, so $r$ is not needed as an accuracy knob; it
is fixed at the minimum that keeps the interaction lists non-degenerate.

\section{Justification of every parameter and design choice}
\label{sec:params}

\subsection{Chebyshev order $q$ (default $6$)}
$q$ sets the rank $r=q^2$ and the accuracy. From \eqref{eq:cheb-bound} accuracy
improves geometrically, while M2M/M2L cost grows like $q^4$ and buffers like
$q^2$. The default $q=6$ ($r=36$) sits at the knee: $\sim10^{-6}$ matvec accuracy
at a few MiB. $q=4$ ($r=16$) is the fast/low-memory setting at $\sim10^{-4}$;
$q\ge8$ buys little for this kernel because the PCG tolerance and the discretisation
error dominate well before the interpolation error.

\subsection{Leaf side $\ell$ (default $8$, $64$ DOF/leaf)}
\label{sec:params-leaf}
$\ell$ trades near against far work. Larger $\ell$ shrinks the tree (fewer boxes,
smaller buffers, $\propto 1/\ell^2$) but inflates the dense near blocks
($\propto\ell^4$ per leaf, $\propto\ell^2$ per DOF). The choice $\ell=8$ keeps
the near matvec ($\propto(2r+1)^2\ell^2 = 9\cdot64$) comparable to the far cost
and gives the $\approx12$ bytes/DOF asymptote with $q=6$. It must be a power of
two dividing $N_s$ (\code{uniform\_quadtree.cpp:16--19}).

\subsection{Near radius $r$ / admissibility $\eta$ (default $r=1$)}
\label{sec:params-eta}
$r=1$ is the smallest radius for which the interaction list is well-defined and
admissible boxes are separated by at least one box width, corresponding to the
$\Hmat$ admissibility parameter $\eta\approx 2$ (diameter $\le \eta\times$
distance). Smaller is impossible; larger only shifts work from far to near at no
accuracy benefit given the Love far kernel (Section~\ref{sec:accuracy}). $r=1$ is
the cost-optimal admissible choice.

\subsection{Geometric box extent}
Using the full geometric side (not the centre bounding box) for normalisation
yields the exact $1/2$ child-to-parent scaling (Proposition~\ref{prop:halfscale}),
which is what collapses M2M/L2L to four level-independent matrices
(Proposition~\ref{prop:m2m-cache}). A centre-based extent would make the scaling
$\ell$- and level-dependent and destroy the cache.

\subsection{bbFMM versus analytic FMM versus ACA}
\emph{bbFMM} needs only kernel \emph{evaluations}, so it accommodates the Love
cell integral (with its logarithms) directly -- no multipole expansion of the
integrated kernel is required, unlike an \emph{analytic FMM}, which would demand
a tailored series for $g$ and re-derivation for any kernel change. Compared with
a classical \emph{ACA} $\Hmat$-matrix, bbFMM gives nested, shared bases and
hence $\Order(N)$ (not $\Order(N\log N)$) memory and matvec, and -- via
translation invariance -- avoids storing redundant per-block factors
(Section~\ref{sec:cache}). It is also kernel-independent at the level of
software: the same machinery applies to any smooth $g$ supplied as a function,
in the spirit of kernel-independent FMM~\cite{Ying2004,FongDarve2009}.

\subsection{Love far kernel rather than $1/r$}
Consistency between near and far (Section~\ref{sec:lovefar}): the far error is
then purely interpolation, with no near/far model mismatch.

\subsection{Additive backend}
The $\Htwo$ operator is an additive backend: it implements the same
matvec interface (\code{H2Operator::matvec}) as the dense and $\Hmat$ paths, so
the PCG solver and all existing tests are untouched and the three backends are
directly comparable on identical problems.

\section{Validation and comparisons}
\label{sec:valid}

\subsection{Against the dense operator}
At small $N$ the apply is checked entrywise against \code{assemble\_dense}:
relative $L_2$ matvec error $1.3\times10^{-4}$ ($q=4$) and $3.2\times10^{-6}$
($q=6$) at $N_s=32$ (\code{tests/test\_h2.cpp}). This certifies the operator
implementation independently of the contact solve.

\subsection{Against the classical $\Hmat$-matrix}
The $\Htwo$ backend reproduces the $\Hmat$-matrix contact area exactly on the
rough-surface benchmark, at a fraction of the memory: where the classical
$\Hmat$ needs $6.2$\,GiB at $N_s=512$, the $\Htwo$ operator needs $5.26$\,MiB
(Table~\ref{tab:mem}) -- three orders of magnitude less -- because it stores
shared bases and cached couplings instead of independent per-block factors.

\subsection{Against the FFT solver (Tamaas)}
For a fully active rectangular grid the operator is a convolution and a
zero-padded FFT matvec is the natural reference. Tamaas~2.8.1 implements this via
its \code{dcfft} operator. Table~\ref{tab:tamaas} compares the $\Htwo$ backend
($q=6$) against Tamaas on identical self-affine surfaces
(\code{compare\_tamaas\_h2.py}).

\begin{table}[h]
\centering
\begin{tabular}{rrrrr}
\toprule
$N_s$ & frac.\ (Tamaas) & frac.\ ($\Htwo$) & rel.\ $L_2$ pressure & $\Htwo$ mem [MiB] \\
\midrule
64  & 0.1265 & 0.1260 & 3.26\% & 1.12 \\
128 & 0.1215 & 0.1201 & 3.17\% & 1.66 \\
256 & 0.1112 & 0.1098 & 2.92\% & 2.61 \\
\bottomrule
\end{tabular}
\caption{$\Htwo$ vs Tamaas FFT. Contact fractions agree to $\sim10^{-3}$
(reproduced from the concise note).}
\label{tab:tamaas}
\end{table}

The residual $\sim3\%$ pressure difference is \emph{not} an $\Htwo$ error: the
$\Htwo$ operator reproduces the exact Love operator to $3\times10^{-6}$ at $q=6$,
whereas Tamaas~2.8.1's \code{dcfft} influence coefficients oscillate $\pm2$--$8\%$
from the exact Love values at $1$--$3$ cell separations (Gibbs-like), and its
effective modulus is $2\Es{}^2$ (both verified against Hertz theory; see the
project README and \code{doc/theory/h2\_fmm.tex}). The difference is therefore
dominated by the FFT reference, and it is the same $\approx3.3\%$ seen for the
$\Hmat$-matrix backend.

Single-thread timing (Table~\ref{tab:tamaas-time}) shows the $\Htwo$ per-matvec
cost is already below the FFT at every size, with the ratio improving from $0.78$
at $N_s=64$ to $0.69$ at $N_s=512$ as the $\Order(N)$ apply overtakes the FFT's
$\Order(N\log N)$. End-to-end the $\Htwo$ solve is $\sim1.5\times$ faster here,
reflecting also a somewhat lower projected-CG iteration count.

\begin{table}[h]
\centering
\begin{tabular}{rrrrrrr}
\toprule
& \multicolumn{3}{c}{per-matvec [ms]} & \multicolumn{3}{c}{end-to-end solve [s]} \\
\cmidrule(lr){2-4}\cmidrule(lr){5-7}
$N_s$ & Tamaas & $\Htwo$ & ratio & Tamaas & $\Htwo$ & ratio \\
\midrule
64  &  0.36 &  0.28 & 0.78 &  0.029 & 0.017 & 0.58 \\
128 &  1.63 &  1.29 & 0.79 &  0.204 & 0.128 & 0.63 \\
256 &  6.92 &  5.59 & 0.81 &  1.230 & 0.870 & 0.71 \\
512 & 37.30 & 25.89 & 0.69 & 10.319 & 6.822 & 0.66 \\
\bottomrule
\end{tabular}
\caption{Single-thread timing, $\Htwo$ ($q=6$) vs Tamaas \code{dcfft}
(\code{compare\_tamaas\_h2.py}), reproduced from the concise note. Per-matvec is
the same operation $\mathbf u=S\pp$ on both sides.}
\label{tab:tamaas-time}
\end{table}

Both $\Htwo$ and FFT are $\Order(N)$ in memory; $\Htwo$ is additionally
\emph{local}, leaving room to skip work outside the contact set (active-set
masking), which a global FFT cannot do without re-padding.

\section{Convergence acceleration: spectral preconditioning and nested-grid continuation}
\label{sec:accel}

The previous sections make a \emph{single} matvec $\mathbf u=S\pp$ optimal:
$\Order(N)$ work and memory. The cost of a contact \emph{solve}, however, is
(iterations)$\times$(matvec), and Table~\ref{tab:cputime} shows the
projected-CG iteration count climbing with $N_s$ -- from $90$ at $N_s=128$ to
$906$ at $N_s=1024$. Beyond $N_s\sim10^3$ this iteration growth, not the
operator, dominates the wall time. This section analyses the growth and develops
two complementary, \emph{optional} accelerations: a spectral (Fourier)
preconditioner and a nested-grid (cascadic / full-multigrid) continuation with
warm starts. Both are exact in the sense that they change only the path, not the
fixed point: with the preconditioner disabled the solver reproduces the original
iteration bit-for-bit (Proposition~\ref{prop:precond-none}).

\subsection{Why the iteration count grows: conditioning of $S$}
\label{sec:cond}

The relevant spectral fact was already noted in Section~\ref{sec:physics}: on the
infinite plane the convolution operator has symbol
$\widehat G(\mathbf q)=2/(\Es\abs{\mathbf q})$, decaying like $1/\abs{\mathbf q}$.
The discrete operator inherits this through its circulant (periodic) part.

\begin{proposition}[Condition number of the discrete operator]
\label{prop:cond}
Let $S_c$ be the circulant (periodic) discretisation of the half-space kernel on
the $N_s\times N_s$ grid, restricted to the subspace orthogonal to the constant
mode. Its spectral condition number satisfies
\begin{equation}
  \kappa_2(S_c) \;=\; \frac{\lambda_{\max}}{\lambda_{\min}}
  \;=\; \frac{\widehat S_c(\mathbf q_{\min})}{\widehat S_c(\mathbf q_{\max})}
  \;=\; \Theta(N_s).
  \label{eq:kappa}
\end{equation}
\end{proposition}
\begin{proof}
$S_c$ is diagonalised by the 2-D DFT; its eigenvalues are the discrete symbol
$\widehat S_c(\mathbf q)$, the FFT of the influence coefficients, which agree with
the continuous symbol $2/(\Es\abs{\mathbf q})$ up to $\Order(1)$ factors on the
resolved band. The eigenvalues are positive and \emph{decreasing} in
$\abs{\mathbf q}$. The largest acts on the lowest nonzero wavenumber
$\abs{\mathbf q}_{\min}=2\pi/L$ (the constant mode $\mathbf q=0$ is excluded by the
mean constraint), the smallest on the Nyquist wavenumber
$\abs{\mathbf q}_{\max}=\sqrt2\,\pi N_s/L$. Hence
$\kappa_2 = \widehat S_c(\mathbf q_{\min})/\widehat S_c(\mathbf q_{\max})
= \abs{\mathbf q}_{\max}/\abs{\mathbf q}_{\min} = \tfrac{\sqrt2}{2}N_s=\Theta(N_s)$.
\end{proof}

\begin{corollary}[CG iteration growth]
\label{cor:cg}
The conjugate-gradient method applied to a system with condition number $\kappa$
reduces the energy-norm error by a factor $\varepsilon$ in at most
$\big\lceil \tfrac12\sqrt{\kappa}\,\ln(2/\varepsilon)\big\rceil$ iterations. With
$\kappa=\Theta(N_s)$ this is $\Order(\sqrt{N_s})$ iterations -- intrinsic to the
$1/\abs{\mathbf q}$ kernel, independent of the matvec implementation.
\end{corollary}

Two caveats sharpen the picture for \emph{contact}. First, the active constraint
turns the linear solve into a sequence of reduced systems on the (changing)
contact set; \emph{active-set discovery} adds iterations on top of
Corollary~\ref{cor:cg}. Second, the empirical growth depends on how the test
problem is posed. The steep $90\!\to\!906$ growth in Table~\ref{tab:cputime} uses
a surface whose spectrum \emph{extends} with $N_s$ (a rougher problem at each
size), so it conflates conditioning with a genuinely harder problem. Isolating the
conditioning term requires a \emph{fixed physical} spectrum: the band-limited
self-affine surface of \code{bandlimited\_surface} (\code{experiments/precond\_pcg.py}),
whose wavenumber band $k\in[4,32]$ is independent of $N_s$, so coarse and fine
grids resolve the \emph{same} surface. On that family the unpreconditioned count
grows as $26\!\to\!43\!\to\!56\!\to\!83\!\to\!180$ for $N_s=64,\dots,1024$
(Table~\ref{tab:precond-it}) -- close to the $\sqrt{N_s}$ of
Corollary~\ref{cor:cg} (a $16\times$ range in $N_s$ giving roughly a $7\times$
range in iterations, the remainder being active-set work). The smooth Hertz
problem grows only mildly with $N_s$, a further confirmation that the dominant
term is the operator conditioning, not the contact logic.

\subsection{The spectral preconditioner}
\label{sec:precond}

\paragraph{Ideal symbol.}
Because $S_c$ is diagonalised by the DFT, the ideal (constant-spectrum)
preconditioner is the operator with the \emph{inverse} symbol,
\begin{equation}
  \widehat{M^{-1}}(\mathbf q) = \frac{1}{\widehat S_c(\mathbf q)},
  \qquad \mathbf q\ne 0,
  \label{eq:Minv-symbol}
\end{equation}
for then $M^{-1}S_c$ has every nonzero eigenvalue equal to $1$ and
$\kappa(M^{-1}S_c)=1$: CG converges in $\Order(1)$ iterations. Two realisations of
\eqref{eq:Minv-symbol} are used.
\begin{itemize}
\item \textbf{Circulant symbol} $w_c(\mathbf q)=1/\widehat S_c(\mathbf q)$, with
$\widehat S_c$ the FFT of the kernel's influence coefficients -- obtainable from
the lookup table $T$ directly, or from a single impulse matvec
$S\boldsymbol\delta$. This is exact for the periodic part of the operator
(\code{symbol\_circulant}, \code{experiments/precond\_pcg.py:66--80}).
\item \textbf{Asymptotic symbol} $w(\mathbf q)=\abs{\mathbf q}$, the continuous
inverse of the $1/\abs{\mathbf q}$ symbol (\code{symbol\_absq}, \code{:57--63};
the integer wavenumber form $\abs{\mathbf k}$ in
\code{fourier\_precond.cpp:12--17}). It is parameter-free and guaranteed
positive.
\end{itemize}
The two coincide for this kernel because $\widehat S_c(\mathbf q)\sim
2/(\Es\abs{\mathbf q})$, so $1/\widehat S_c\sim\tfrac{\Es}{2}\abs{\mathbf q}$:
both have the same shape and differ only by a constant factor, which is
immaterial (Proposition~\ref{prop:scale}). The measured iteration counts of the
two symbols are within one iteration at every grid size
(Table~\ref{tab:precond-it}), so the simpler, always-positive $\abs{\mathbf q}$
form is the production choice and the C++ implementation
(\code{FourierPreconditioner}) carries only it.

\begin{proposition}[Scale invariance]
\label{prop:scale}
Replacing the preconditioner $M^{-1}$ by $cM^{-1}$ for any constant $c>0$ leaves
every iterate $\pp^{(k)}$ of the preconditioned Polonsky--Keer CG unchanged.
\end{proposition}
\begin{proof}
Write $z=M^{-1}g$, so $z\mapsto cz$. The conjugacy coefficient
$\beta=\langle z,g-g_{\text{prev}}\rangle/\langle z_{\text{prev}},g_{\text{prev}}
\rangle$ is a ratio of two $M^{-1}$-linear forms and is therefore homogeneous of
degree $0$ in $c$: $\beta$ is unchanged. The search direction $t=z+\beta t_{\text{prev}}$
scales by $c$ (inductively, $t_{\text{prev}}\mapsto ct_{\text{prev}}$). The step
length $\tau=\langle g,t\rangle/\langle St-\overline{St},t\rangle$ has numerator
$\propto c$ and denominator $\propto c^2$, so $\tau\mapsto\tau/c$, and the update
$\pp\mapsto\pp-\tau t$ uses $\tau t\mapsto(\tau/c)(ct)=\tau t$: invariant. The
projection, overlap correction, and load balance act on $\pp$ alone and are
therefore also invariant.
\end{proof}

\paragraph{The DC mode and the rigid-body constraint.}
The symbol is zeroed at $\mathbf q=0$ ($w(0)=0$; \code{fourier\_precond.cpp:17},
\code{:79}). The mean of the pressure (the total load) and the rigid-body approach
are fixed by the equality constraint $\operatorname{mean}(\pp)=\bar p$ and the
shift $\alpha$ in the gradient, \emph{not} by the CG search; removing the
$\mathbf q=0$ component of the preconditioned residual keeps the CG direction
inside the load-conserving subspace and prevents the (formally infinite,
$1/\widehat S_c(0)$) DC gain from corrupting it.

\paragraph{Application to the contact-masked residual.}
Per iteration the preconditioner forms
\begin{equation}
  z \;=\; \Pi_{\mathcal A}\,\Big(\mathcal F^{-1}\big[\,w(\mathbf q)\,\mathcal F[\,
  \Pi_{\mathcal A} g\,]\,\big]\Big)\;-\;\bar z_{\mathcal A}\,\mathbf 1_{\mathcal A},
  \label{eq:apply}
\end{equation}
where $\mathcal A=\{i:p_i>0\}$ is the current contact (active) set,
$\Pi_{\mathcal A}$ restricts a field to $\mathcal A$ (zero elsewhere),
$\mathcal F$ is the 2-D DFT, and $\bar z_{\mathcal A}$ is the mean of $z$ over
$\mathcal A$. In words: mask the gradient to the contact set, apply the symbol by
FFT, keep the result on the contact set, and subtract its contact-set mean. This
is \code{FourierPreconditioner::apply\_into} (\code{fourier\_precond.cpp}); the
final contact-mean removal is the discrete analogue of the $\mathbf q=0$ zeroing
on the \emph{reduced} space.

Since the July 2026 implementation pass the transform exploits that the masked
residual is \emph{real}: the $x$-direction runs as real-to-complex /
complex-to-real 1-D FFTs, storing only the $n_h=N_s/2+1$ non-redundant $k_x$
modes (the rest follow from Hermitian symmetry), and each stored $k_x$ line is
processed in one fused sweep along $y$ --- forward complex FFT, multiplication
by the real, even symbol $\abs{\mathbf k}$, inverse FFT --- before the
complex-to-real pass returns to the grid. This halves the FFT work and shrinks
the transform scratch from ${\approx}4N$ to ${\approx}2N$ reals ($N_s^2$ real
$+\;n_hN_s$ complex). The scratch is owned by the \code{FourierPreconditioner}
object and reused across CG iterations, and \code{apply\_into} writes into a
caller-owned output buffer, so a preconditioner application performs no heap
allocation in steady state; all passes are OpenMP-parallel with per-thread FFT
plans. The finite-precision aspects (double accumulation of the reductions,
the cancellation-safe line-search denominator) are documented in the companion
reference \code{doc/theory/pcg.tex}, \S``Spectral Preconditioning and
Finite-Precision Implementation''.

\begin{remark}[Why masking is only an approximate diagonaliser]
\label{rem:mask}
The DFT diagonalises the \emph{full periodic} operator $S_c$, not its restriction
$P_{\mathcal A}S P_{\mathcal A}$ to the contact set, which is what the projected
CG actually inverts. Applying the full-grid inverse symbol to the masked residual,
\eqref{eq:apply}, is therefore an \emph{approximate} preconditioner of the reduced
Hessian: it is exact in the interior of a large, simply connected contact patch
(where the field is locally periodic) but leaks at the contact boundary and for
fragmented contact, where masking and the global FFT do not commute. Consequently
$\kappa(M^{-1}P_{\mathcal A}SP_{\mathcal A})$ is reduced to $\Order(1)$ but not to
exactly $1$; the residual growth in Table~\ref{tab:precond-it} (preconditioned
count still rising slowly with $N_s$) is the quantitative footprint of this
boundary leakage. This is the fundamental reason a pure spectral preconditioner
\emph{cannot} be fully grid-independent for a constrained problem
(Section~\ref{sec:accel-limit}).
\end{remark}

\subsection{The minimal change to the Polonsky--Keer PCG}
\label{sec:pcg-change}

The preconditioner enters the solver of Section~\ref{sec:discretisation}
(\code{contact\_solver.cpp}) through a single substitution: the search direction
is built from the preconditioned residual $z$ instead of the bare gradient $g$,
with the conjugacy taken in the $M$-inner product. Everything else -- in
particular the \emph{exact} line search, which is what makes Polonsky--Keer robust
-- is untouched. Per iteration, on the contact set $\mathcal A$:
\begin{align}
  g &= S\pp + \mathbf g_0 - \alpha, \qquad
     \alpha=\operatorname{mean}_{\mathcal A}(S\pp+\mathbf g_0)
     &&\text{(gradient, rigid shift; unchanged)}\notag\\
  z &= M^{-1} g \quad\text{by \eqref{eq:apply}}
     &&\text{(\textbf{new}; $z=g$ on $\mathcal A$ if no precond)}\notag\\
  \beta &= \max\!\Big(0,\ \frac{\langle z,\,g-g_{\text{prev}}\rangle_{\mathcal A}}
                              {\langle z_{\text{prev}},\,g_{\text{prev}}\rangle_{\mathcal A}}\Big)
     &&\text{(\textbf{modified}: PR$^+$ in the $M$-inner product)}\notag\\
  t &= z + \beta\, t_{\text{prev}} \ \ \text{on }\mathcal A,\quad 0\text{ off }\mathcal A
     &&\text{(\textbf{modified}: was $g+\beta t$)}\notag\\
  \tau &= \frac{\langle g,t\rangle_{\mathcal A}}
              {\langle St-\overline{St},\,t\rangle_{\mathcal A}}
     &&\text{(\textbf{unchanged}: exact line minimiser)}\notag\\
  \pp &\leftarrow \big(\pp-\tau t\big)_+,\quad
       p_i \mathrel{-}=\tau g_i\ \text{for released penetrating }i,\quad
       \pp \mathrel{*}= \tfrac{P_{\text{tot}}}{\textstyle\sum p}
     &&\text{(\textbf{unchanged})}\notag
\end{align}
This is exactly \code{contact\_solver.cpp:63--122}: the preconditioned residual
$z$ is computed at \code{:64--69} (with the unpreconditioned branch
\code{z(i)=contact[i]?g(i):0} reproducing $z=g$ on $\mathcal A$); the $M$-inner
products $G=\langle z,g\rangle_{\mathcal A}$ and
$G_{\text{pr}}=\langle z,g-g_{\text{prev}}\rangle_{\mathcal A}$ at \code{:72--78};
the direction $t=z+\beta t$ at \code{:80--81}; and the unchanged line search at
\code{:86--101}.

\begin{proposition}[Exact reduction to the original solver]
\label{prop:precond-none}
With \code{precond} empty (\code{precond="none"}), the preconditioned solver
executes the \emph{same algorithm} as the original Polonsky--Keer PCG: every
iterate coincides in exact arithmetic.
\end{proposition}
\begin{proof}
With no preconditioner the code sets $z_i=g_i$ on $\mathcal A$ and $0$ elsewhere.
Then $G=\langle g,g\rangle_{\mathcal A}$ and
$G_{\text{pr}}=\langle g,g-g_{\text{prev}}\rangle_{\mathcal A}$ are precisely the
Fletcher--Reeves / Polak--Ribière$^+$ scalars of the unmodified algorithm, and
$t=g+\beta t$ is the original direction. No other step is touched. (Verified
numerically in \code{tests/test\_precond.cpp}, and as the
\code{cpp}$=$\code{none} columns of Table~\ref{tab:precond-it}.)
\end{proof}

\begin{remark}[Bit-reproducibility]
Until June 2026 the reduction was also bit-for-bit: the unpreconditioned branch
reproduced the historical solver's floating-point operations one by one. The
July 2026 performance pass parallelised all $O(N)$ vector passes with OpenMP
and accumulates the scalar reductions in \code{double} independent of the
working precision, so the summation \emph{order} now depends on the thread
count. The iterates therefore agree with the historical solver to rounding
level rather than bitwise (measured pressure difference
$9\times10^{-15}$ relative $L^2$ at $N_s=1024$, identical contact set); in
single precision the double accumulation is a strict accuracy \emph{gain}
(see \code{doc/theory/pcg.tex}).
\end{remark}

\subsection{Warm start and nested-grid (cascadic / FMG) continuation}
\label{sec:nested}

The preconditioner attacks conditioning; the nested-grid continuation attacks the
\emph{remaining} $N_s$-dependence (active-set discovery and boundary leakage,
Remark~\ref{rem:mask}) by never solving a hard fine problem cold. The idea is the
full-multigrid / cascadic strategy~\cite{Brandt1977,BornemannDeuflhard1996}: solve
the cheapest (coarsest) grid first and use its solution to warm-start the next
finer grid, sweeping coarse$\to$fine once.

\paragraph{Grid transfers.}
Let level $\ell$ have grid size $N_s^{(\ell)}=2^{\ell}N_s^{(0)}$.
\begin{itemize}
\item \textbf{Restriction} (fine$\to$coarse), applied to the \emph{surface} so
that every level is the same physical problem: the $2\times2$ block average
$\big(R h\big)_{IJ}=\tfrac14\sum_{a,b\in\{0,1\}}h_{2I+a,\,2J+b}$
(\code{restrict}, \code{experiments/nested\_grid.py:28--32}). The finest surface
is generated once and restricted down to each coarse level (\code{:56--60}), so
coarse and fine grids see identical low-frequency content.
\item \textbf{Prolongation} (coarse$\to$fine) of the converged \emph{pressure},
by \emph{injection}: each coarse value is replicated into its $2\times2$ fine
block, $\big(P\pp_c\big)_{2I+a,\,2J+b}=(\pp_c)_{IJ}$
(\code{prolong(..., mode="inject")}, \code{:35--41}).
\end{itemize}

\begin{lemma}[Injection is load-preserving and active-set-sharp]
\label{lem:inject}
Injection prolongation preserves the total load and maps the coarse contact set
onto its exact geometric refinement.
\end{lemma}
\begin{proof}
A coarse cell of area $h_c^2$ refines into four fine cells of area
$h_f^2=(h_c/2)^2=h_c^2/4$. Injecting the coarse pressure $p$ into all four gives a
fine contribution $4\cdot p\cdot h_f^2 = p\,h_c^2$, equal to the coarse cell's
load; summing over cells, the total load is preserved. A fine cell is active
($p>0$) iff its parent coarse cell is active, so the prolonged active set is the
pixel-exact refinement of the coarse contact patch -- in particular its boundary
stays sharp.
\end{proof}

\begin{remark}[Why not bilinear]
Bilinear prolongation produces fine pressures that are convex combinations of
neighbouring coarse values. Across the contact boundary it manufactures small
\emph{positive} pressures in cells whose true state is open (and shaves the
boundary cells), so the warm-started active set is blurred. Since the active set,
not the pressure magnitude, is the hard latent variable of the contact problem, a
blurred seed costs the fine solve extra iterations to re-sharpen the boundary.
Lemma~\ref{lem:inject} is exactly the property bilinear lacks. The measurements
confirm the prediction at every size: injection beats bilinear ($26$ vs $30$ at
$N_s=256$, $31$ vs $32$ at $512$, $45$ vs $51$ at $1024$;
Table~\ref{tab:precond-it}).
\end{remark}

\paragraph{Warm start and cascadic tolerance.}
At each level the solver is launched from \code{p\_init} $=$ the prolonged coarse
pressure (renormalised to the target load, \code{contact\_solver.cpp:21--27}),
seeding both the pressure field and the active set; \code{p\_init = nullptr}
recovers the cold uniform start $p\equiv\bar p$. Every level uses the
$\abs{\mathbf q}$ preconditioner. Coarse levels are solved only to a loose
tolerance ($10^{-4}$ in the prototype, \code{experiments/nested\_grid.py:68}), the
\emph{cascadic} choice~\cite{BornemannDeuflhard1996}: a coarse solution accurate
enough to seed the next level need not be converged to the final tolerance, which
is applied only on the finest grid. Finally, the fine-level overlap correction
(\code{contact\_solver.cpp:108--115}) can still \emph{reactivate} nodes the coarse
seed left open, so the warm start is a seed, not a hard constraint: convergence to
the correct fine fixed point is unaffected.

\begin{remark}[Single-entry implementation]
\label{rem:nested-api}
The entire coarse$\to$fine continuation is implemented as one C++ entry point,
with no Python orchestration: \code{hmc::solve\_contact\_nested(...)}
(\code{include/nested\_solve.hpp}), exposed to Python as
\code{hc.solve\_nested(grid\_size, gap, p\_nominal, coarsest=64, q=6, ...)}.
It builds the grid hierarchy and the per-level $\Htwo$ operators internally,
restricts the gap by $2\times2$ averaging, and warm-starts each level by
injecting the prolonged coarse pressure. It requires
$\texttt{grid\_size}=\texttt{coarsest}\cdot2^{k}$ and returns a standard
\code{ContactResult}.
\end{remark}

\subsection{Results}
\label{sec:accel-results}

\paragraph{Prototype (Python), fixed-band rough surface.}
With $p_{\text{nominal}}=0.05$, tol $10^{-8}$, $q=6$
(\code{experiments/precond\_pcg.py}, \code{nested\_grid.py};
\code{experiments/precond\_results.md}):

\begin{table}[h]
\centering
\begin{tabular}{rrrrrrrr}
\toprule
$N_s$ & baseline & $\abs{q}$ & circulant & nested (inj.) & nested (bilin.)
& $A_c/A$ & rel-$p$ \\
\midrule
64   & 26  & 19 & 19 & ---  & ---  & 0.125 & $1.4\times10^{-7}$ \\
128  & 43  & 25 & 24 & ---  & ---  & 0.098 & $2.4\times10^{-7}$ \\
256  & 56  & 33 & 33 & 26   & 30   & 0.091 & $3.3\times10^{-7}$ \\
512  & 83  & 43 & 43 & 31   & 32   & 0.095 & $4.5\times10^{-7}$ \\
1024 & 180 & 62 & 62 & 45   & 51   & 0.105 & $4.9\times10^{-7}$ \\
\bottomrule
\end{tabular}
\caption{Projected-CG iteration counts on the fixed-band rough surface.
``baseline'' is the unpreconditioned solver (which equals the C++ \code{none}
count exactly, Prop.~\ref{prop:precond-none}); $\abs{q}$ and ``circulant'' are
the two preconditioner symbols; ``nested'' is the \emph{fine-grid} iteration count
of the cascadic continuation (coarsest $N_s{=}64$) with injection vs bilinear
prolongation. $A_c/A$ is the contact fraction and rel-$p$ the pressure rel-$L_2$
of the preconditioned vs the unpreconditioned solution -- confirming the
acceleration does not change the result.}
\label{tab:precond-it}
\end{table}

Three readings of Table~\ref{tab:precond-it}.
(i)~\emph{Preconditioning halves the growth rate}: the speedup rises from
$1.37\times$ ($N_s=64$) to $2.90\times$ ($N_s=1024$), because the preconditioned
count grows roughly half as fast as the baseline -- the advantage compounds at
scale, which is exactly the large-grid regime that motivated it.
(ii)~\emph{The warm start cuts a further $\sim1.4\times$} on top of the
preconditioner ($62\!\to\!45$ at $N_s=1024$); combined with preconditioning this
is $\sim4\times$ fewer fine iterations than the original cold/unpreconditioned
baseline ($180\!\to\!45$).
(iii)~The \emph{whole} coarse$\to$fine solve costs \emph{less} than a single cold
solve on the target grid: the cost-weighted work ratio
$\sum_\ell\text{(iters}_\ell)\,N_\ell^2 \,/\, \text{(cold iters)}\,N_{\text{target}}^2$
is $0.89\times$ at $N_s=256$ and $0.79\times$ at $N_s\ge512$
(\code{nested\_grid.py:98--100}). Multigrid pays for itself: the coarse levels are
geometrically cheaper ($N_\ell^2$ shrinks by $4$ per level) and the seed they
provide more than repays their cost.

\paragraph{C++ port.}
The $\abs{\mathbf q}$ preconditioner and the warm-start path are implemented in
the production solver (\code{fourier\_precond.\{hpp,cpp\}}; the \code{precond} and
\code{p\_init} arguments of \code{solve\_contact}, \code{contact\_solver.hpp:38--42}).
The port matches the prototype iteration-for-iteration and is
\emph{solution-identical} to the unpreconditioned solver: contact-area change $0$
and pressure rel-$L_2\sim5\times10^{-7}$ (\code{experiments/bench\_cpp\_precond.py};
the regression test \code{tests/test\_precond.cpp} asserts rel-$p<10^{-5}$,
$\abs{\Delta A_c}<10^{-3}$, and that the preconditioner never increases the
iteration count). In wall-clock terms at $N_s=1024$ on the fixed-band surface the
solve drops from $10.2$\,s (\code{none}) to $7.9$\,s (\code{fourier}) to $6.5$\,s
(nested), an end-to-end $\sim1.6\times$ -- the iteration-count gain translated into
runtime through the $\Order(N)$ matvec. On the smooth Hertz problem at $N_s=64$
the preconditioner alone cuts $27\!\to\!16$ iterations, and warm-starting from the
converged pressure converges in $0$ iterations (the initial residual is already
below tolerance), the trivial limit of the continuation.

\subsection{Honest limitation and outlook for the acceleration}
\label{sec:accel-limit}

Neither device makes the solve fully grid-independent, and for a principled
reason. The spectral preconditioner is limited by the contact restriction
(Remark~\ref{rem:mask}): masking the residual to the contact set breaks the FFT
diagonalisation at the contact boundary and for fragmented contact, so the
preconditioned reduced system is well-conditioned but not perfectly so. The
nested continuation is limited by the physics: each octave of refinement resolves
genuinely \emph{new} asperities, so there is real new fine-scale work at every
level -- the fine-grid count still rises ($26\!\to\!45$ over $N_s=256\!\to\!1024$,
Table~\ref{tab:precond-it}). This residual growth is inherent to rough contact,
not a defect of the method.

These two accelerations are, however, the practical path \emph{toward} near
grid-independence, banking a measured $\sim4\times$ in fine iterations and
$\sim1.6\times$ in wall time at $N_s=1024$ at zero change to the solution and at
the cost of one FFT per iteration. The heavier alternative is a genuine
\emph{monotone / constraint multigrid} for the variational
inequality~\cite{Kornhuber1994}, which builds the inequality constraints into the
coarse-grid corrections (truncated/monotone restrictions) rather than only seeding
the active set; that is expected to remove the remaining boundary-leakage term at
the cost of a substantially more intricate solver, and is left as future work.

\section{Limitations and outlook}
\label{sec:outlook}

The operator currently runs on the full grid; the contact constraint is handled
by the PCG projection, not by skipping inactive DOFs. Natural extensions, none of
which change the public interface
(\code{H2Operator\{build(),matvec(),print\_statistics()\}}):
\begin{itemize}
\item \textbf{Active-set masking.} Restrict the near and far passes to leaves and
interactions touching the contact patch. Because the $\Htwo$ apply is local
(unlike the FFT), masking is natural and can reduce the per-iteration cost as the
contact area shrinks.
\item \textbf{Rectangular / refined grids.} The uniform quadtree currently
requires a square, power-of-two grid with $\ell\mid N_s$
(\code{uniform\_quadtree.cpp:16--19}). Anisotropic ($n_x\ne n_y$) and locally
refined trees would broaden applicability at the cost of more distinct cached
couplings.
\item \textbf{Single-precision caches.} Storing couplings, stencils, and buffers
in \code{float} halves memory; given the $10^{-6}$ accuracy target at $q=6$,
single precision in the far field is comfortably adequate.
\item \textbf{FFT backend for the fully-active case.} For a saturated rectangular
contact, the DC-FFT convolution remains competitive and could be offered
alongside, with the $\Htwo$ operator preferred for sparse contact and the FFT for
near-full contact.
\end{itemize}

\section*{Summary}
The Hcontact $\Htwo$/FMM operator applies the Boussinesq flexibility matrix in
$\Order(N)$ time and memory, matrix-free, by combining (i) an exact short-range
near field from the Love cell-integral lookup table; (ii) a tensor-product
Chebyshev (bbFMM) far field with spectral accuracy in the order $q$; and (iii) a
translation-invariant caching that collapses the transfer and coupling operators
to four matrices, $\Order(\log N)$ couplings, and $\Order(1)$ near stencils. The
near/far split tiles every pair of unknowns exactly once
(Theorem~\ref{thm:coverage}), so the only approximation is a controlled,
$N$-independent interpolation error ($3.2\times10^{-6}$ at $q=6$). The operator
reaches $2.7\times10^8$ unknowns at second-scale build and matvec times and
$14$\,GiB, regimes entirely out of reach for dense or classical $\Hmat$
representations, and is faster per matvec than a state-of-the-art FFT reference.

\bibliographystyle{unsrt}
\bibliography{references}

\end{document}
